\section{Reflection positivity and the plane covariance}
\label{sec:spectral}

All expectations in this section are for the free plane pair constructed in
\cref{prop:plane-mixing}. Its height increments are denoted by
\(\xi_u=H(u)-H(0)\), for \(u\in\Tri\). For a compactly supported test of integral
zero, \(H_\delta(f)\) means the integral of the affine interpolation on
\(\delta\Tri\) against \(f\); its value is independent of the chosen additive
constant. No Gaussian assumption is made in this section.

We identify the scaling covariance of these averages in two stages. First,
reflection positivity and an angular inequality force every subsequential
spectral limit to have an inverse-Laplacian density. A signed angular sum
then shows that all these limits have the same coefficient. The covariance
also implies that \(H_\delta(\Delta\phi)\), for a compactly supported smooth
\(\phi\) supported strictly on one side of a lattice reflection line, has
vanishing reflection norm.
This is the input for the pin-motion argument in \cref{sec:pins}.

We use Euclidean coordinates in both spaces. Set
\[
 e_1=(1,0),\qquad e_2=(1/2,\sqrt3/2),\qquad
 \Tri^\perp=\{p\in\R^2:p\cdot u\in2\pi\Z\text{ for all }u\in\Tri\},
 \qquad \widehat\Tri=\R^2/\Tri^\perp.
\]
The characters are \(u\mapsto e^{ip\cdot u}\). A neighborhood of zero in
\(\widehat\Tri\) will always be identified with its Euclidean representative.
In particular, the disk of radius \(1/2\) is unambiguous. Our continuum Fourier
convention is
\[
 \widehat f(\zeta)=\int_{\R^2}e^{-i\zeta\cdot x}f(x)\,dx.
\]

\subsection{A positive normal translation}

The square-integral proof of reflection positivity below is a site-plane
version of the classical mechanism developed by
Fr\"ohlich--Israel--Lieb--Simon \cite{FILS1978}.  We give the finite-volume
factorization and the exact row geometry because the later argument needs
a positive normal transfer, including in the presence of auxiliary
height phases.
The distinction between site reflection and reflection between site rows
matters: the positive-transfer construction in
\cite[Sections~3.3 and 4.1]{Usui2012} makes the role of site reflection
explicit.  Our halfway line is itself a row of triangular-lattice sites.

\begin{proposition}[Reflection positivity with auxiliary phases]
\label{prop:reflection-positivity}
Let \(\ell\) be a line of lattice sites parallel to a nearest-neighbor direction.
The plane law is reflection positive across \(\ell\). This remains true after
adjoining any finite collection of phase fields
\begin{equation}
 \label{eq:auxiliary-phases}
 Z_\theta(u)=U_\theta\exp\{i\theta(H(u)-H(o))\},\qquad \theta\in\R,
\end{equation}
where the \(U_\theta\)'s are independent uniform points of the unit circle,
independent of the pair, and \(o\) is any anchor.

More precisely, let \(\Theta_\ell\) denote spatial reflection, and let
\(\mathcal A_\ell\) consist of cylinder functions supported in one closed
half-plane, with increments evaluated along paths in that half-plane. Then
\[
 \langle F,G\rangle_\ell
   =\EE\big[\overline{\Theta_\ell F}\,G\big]
\]
is positive semidefinite. In its quotient completion, normal translation farther
into the half-plane by \(\sqrt3\) is a positive self-adjoint contraction \(S\).
Tangential translation by one lattice unit is a commuting unitary operator.
On mesh \(\delta\), the normal step is \(\sqrt3\delta\).
\end{proposition}

\begin{proof}
The phase law does not depend on its anchor: replacing \(o\) by another site
multiplies each \(U_\theta\) by a unit complex number determined by the
increments, which preserves its conditional Haar law. The same observation
under translations and reflections gives the phase law the spatial
symmetries of the pair law. We may therefore put the phase anchor on
\(\ell\).

First take a finite reflection-symmetric free volume. Conditional on the spin
values on its intersection with \(\ell\), the configurations on the two sides
are independent and reflected-identical. Indeed, no nearest-neighbor edge
jumps across a row of sites. A lift on the row is fixed by its spin values and
one anchor value on the row. If phases are present, additionally condition on
their values at that anchor. The phases on each side are then determined by
the increments on that side and these row values, so conditional independence
persists. Thus
\[
 \EE\big[\overline{\Theta_\ell F}F\big]
   =\EE\big[\lvert\EE[F\mid\text{row data}]\rvert^2\big]\ge0.
\]
The plane limit gives the assertion for cylinder functions. Functions in
ordinary \(L^2\) obtained by approximation are included because
\(\lvert\langle F,G\rangle_\ell\rvert\le
\|F\|_{L^2}\|G\|_{L^2}\).

For clarity, take \(\ell=\{y=0\}\), with the upper half-plane chosen. Write a
site as \((m+n/2,n\sqrt3/2)\). Reflection about \(y=0\) sends its integer
coordinates to \((m+n,-n)\), and reflection about \(y=\sqrt3/2\) sends them to
\((m+n-1,2-n)\). Both preserve \(\Tri\). Their composition is the pure normal
translation \(a=(0,\sqrt3)\); no tangential correction is required.
Let \(T\) be translation by \(a\) on cylinder functions. Stationarity gives
\(\langle TF,G\rangle_\ell=\langle F,TG\rangle_\ell\). Moreover,
\(\langle F,TF\rangle_\ell\) is the reflection form of \(TF\) at the halfway
row, since reflecting \(TF\) in that row gives \(\Theta_\ell F\). It is
therefore nonnegative.

To prove boundedness before taking the quotient, put
\(a_n=\|T^nF\|_\ell\). Reflection Cauchy--Schwarz and the preceding symmetry
imply \(a_n^2\le a_{n-1}a_{n+1}\) for \(n\ge1\). Ordinary Cauchy--Schwarz
and stationarity give \(a_n\le\|F\|_{L^2}\) for every \(n\). A bounded
nonnegative log-convex sequence cannot have \(a_1>a_0\); if \(a_0=0\), the
same inequality gives \(a_1=0\). Consequently \(\|TF\|_\ell\le\|F\|_\ell\).
It follows that \(T\) descends to the quotient and extends to the stated
positive self-adjoint contraction \(S\). Tangential translation commutes with
reflection, preserves the form, and has an inverse with the same properties.
It is therefore unitary and commutes with \(S\). Rotation gives the other
lattice directions, and rescaling gives the mesh statement.
\end{proof}

\subsection{The increment spectral measure}

We first specify the elementary spectral construction for ordinary lattice
translations. If $V$ is a unitary translation and $F$ a Hilbert-space
vector, the positive Fourier polynomials
\[
 \frac1{N+1}\left\|\sum_{j=0}^N e^{-ijk}V^jF\right\|^2
             \frac{dk}{2\pi}
\]
have mass $\|F\|^2$ and Fourier coefficients converging to
$\langle F,V^nF\rangle$.  Compactness of finite measures on the circle
gives the scalar spectral measure, uniquely specified by these
coefficients.  Rectangular Fourier sums give the same construction for
commuting lattice translations.  Polarization gives their cross
measures and the bound on total variation by the product of the vector
norms.  A multiplier is defined first for trigonometric polynomials and
then by completion; this is the meaning of $E(dp)$ below.  The mass at
zero is the projection onto invariant vectors, also obtained as the
limit of rectangular Ces\`aro averages of translations.

\begin{lemma}
\label{lem:increment-spectral}
There is a nonnegative measure \(C\), locally finite on
\(\widehat\Tri\setminus\{0\}\), such that
\begin{equation}
 \label{eq:increment-spectral}
 \EE\left(\sum_u a_uH(u)\right)^2
   =\int_{\widehat\Tri\setminus\{0\}}
      \left|\sum_u a_ue^{ip\cdot u}\right|^2 C(dp)
\end{equation}
for every finitely supported real array with \(\sum_u a_u=0\).
The measure is invariant under the lattice reflections and rotations.
There is no additional random linear tilt term in this identity.
\end{lemma}

\begin{proof}
Let \(V_u\) be the ordinary unitary translation operators on the plane
probability space, and let \(E(dp)\) be their joint spectral resolution. The
increment cocycle satisfies
\[
 \xi_{u+v}=\xi_v+V_v\xi_u,
 \qquad (V_v-I)\xi_u=(V_u-I)\xi_v.
\]
On a spectral neighborhood where \(e^{ip\cdot v}-1\ne0\), this identity
expresses every \(E(dp)\xi_u\) as \(e^{ip\cdot u}-1\) times the same spectral
vector. Equivalently, on that neighborhood define
\[
 C(dp)=\frac{\langle\xi_v,E(dp)\xi_v\rangle_{L^2}}
                  {|e^{ip\cdot v}-1|^2}.
\]
The cocycle identity makes these definitions agree on overlaps. The two basis
directions cover the punctured torus, so they define a locally finite measure
there and give \cref{eq:increment-spectral} for the nonzero spectral part.

Let \(P_0=E(\{0\})\). The invariant part of the cocycle is additive, hence
\(P_0\xi_u=A\cdot u\) for a random vector \(A\in L^2\). Take a smooth
zero-integral function \(f\) with nonzero first moment \(m=\int x f(x)\,dx\),
and smear the unit-mesh affine height against \(L^{-2}f(x/L)\). Affine
interpolation reproduces linear functions exactly, so the invariant part of
this statistic is \(L A\cdot m\). Its \(L^2\) norm is bounded uniformly in
\(L\) by \cref{prop:plane-moments}, since orthogonal projection cannot
increase that norm. Two choices of \(f\) with independent first moments imply
\(A=0\). This also rules out cancellation of a tilt by other spectral
components: those components are orthogonal to the invariant subspace.

The representation determines \(C\) locally by any nonvanishing lattice
difference. The symmetries of the increment law therefore give the asserted
symmetries of \(C\).
\end{proof}

\begin{lemma}[Uniform annular mass]
\label{lem:annulus-mass}
For some finite constant \(K\),
\begin{equation}
 \label{eq:annulus-mass}
 C\{r\le |p|\le2r\}\le K,\qquad 0<r<1/4.
\end{equation}
In particular, for every \(\eta>0\),
\(\int_{0<|p|<1/4}|p|^\eta C(dp)<\infty\).
\end{lemma}

\begin{proof}
Choose two real smooth zero-integral functions \(f_1,f_2\) whose Fourier
transforms have no common zero on \(1\le|\zeta|\le2\). For example, take the
two first derivatives of a smooth bump whose transform does not vanish on that
annulus. Let \(a^{(r,i)}_u\) be the affine interpolation weights for
\(r^2f_i(rx)\). Uniformly on the fixed annulus,
\[
 \sum_u a^{(r,i)}_ue^{ir\zeta\cdot u}
        \longrightarrow \widehat f_i(-\zeta).
\]
This follows directly by interpolating the smooth plane wave; its error on a
unit triangle is \(O(r^2)\), uniformly for \(\zeta\) in a compact set.
Consequently the sum of the squared absolute values of these two transforms
is bounded below on \(r\le|p|\le2r\), for all sufficiently small \(r\).
Apply \cref{eq:increment-spectral} and the uniform \(L^2\) bound of
\cref{prop:plane-moments}. The remaining range of \(r\) is covered by local
finiteness of \(C\). Summing over dyadic annuli proves the last assertion.
\end{proof}

\subsection{Cauchy kernels from positive transfer}

Transfer spectral representations for lattice two-point functions are
developed in \cite{Usui2012} and, for the six-vertex model, in
\cite[Theorems~4.12 and 4.15]{DCKLM2026}.  We derive the particular
periodized Cauchy mixture needed here, including its reciprocal-lattice
aliases, from the normal step already established.

For the normal transfer we need the joint spectrum of a commuting unitary
$U$ and positive contraction $S$. The preceding Fourier construction
extends to a measure on the circle times $[0,1]$. For each $N$, the commuting
Bernstein operators
\[
 B_{N,j}=\binom Nj S^j(I-S)^{N-j},\qquad 0\le j\le N,
\]
are positive and sum to the identity. For a Hilbert-space vector $F$, place
at $j/N$ the positive Fourier measure whose $n$th coefficient is
\[
                    \langle F,U^nB_{N,j}F\rangle.
\]
The resulting measures have mass $\|F\|^2$. Their mixed moments converge
to $\langle F,U^nS^mF\rangle$, because
$\sum_j(j/N)^m B_{N,j}=S^m+O_m(N^{-1})$ in operator norm.
Compact limits are unique: polynomials in the second coordinate and
trigonometric polynomials in the first uniformly approximate continuous
functions. Polarization again supplies cross measures. For $\Rea z>0$,
integrating $\lambda^z$, with value zero at $\lambda=0$, against them gives
the bounded complex transfer powers used in \cref{sec:pins}. These finite positive-measure
constructions suffice; no reconstruction of a continuum field is used.

Put \(P=2\pi/\sqrt3\). Restrict translations to the rectangular sublattice
generated by \((1,0)\) and \((0,\sqrt3)\), and let \(C'\) be the projection
of \(C\) onto its reciprocal rectangle
\[
 \mathcal R=[-\pi,\pi)\times[-P/2,P/2).
\]
The projection is a two-sheeted covering of reciprocal tori. Its second
preimage of the rectangular origin is a nonzero point of \(\widehat\Tri\).
It will be called an alias of the origin. In every comparison below the
measure near this nonzero point is finite; atoms there are allowed.

For \(0<s<\infty\), let \(Q_s\) be the probability measure on the vertical
circle with density
\begin{equation}
 \label{eq:periodized-cauchy}
 q_s(l)=\sum_{n\in\Z}\frac{s}{\pi\{s^2+(l+nP)^2\}},
       \qquad -P/2\le l<P/2.
\end{equation}
Set \(Q_0=\delta_0\) and let \(Q_\infty\) be uniform probability on that circle.

\begin{proposition}[Positive Cauchy mixture]
\label{prop:cauchy-mixture}
There is a nonnegative measure \(\nu(dk,ds)\) on
\((\R/2\pi\Z\setminus\{0\})\times[0,\infty]\), finite on sets whose
horizontal projection stays away from zero, such that
\begin{equation}
 \label{eq:cauchy-mixture}
 C'(dk,dl)=\int_{[0,\infty]}Q_s(dl)\,\nu(dk,ds),\qquad k\ne0.
\end{equation}
Moreover,
\begin{equation}
 \label{eq:weighted-nu}
 \int_{0<|k|<\kappa}|k|\,\nu(dk,ds)<\infty
\end{equation}
for every sufficiently small fixed \(\kappa>0\).
\end{proposition}

\begin{proof}
Let \(F=H(e_1)-H(0)\), supported on the reflection row. In the reflection
Hilbert space, let \(U\) be horizontal unit translation and \(S\) the positive
normal transfer from \cref{prop:reflection-positivity}. For \(j\ge0\),
\[
 \EE\big[F\,V_{ne_1+j(0,\sqrt3)}F\big]
       =\langle F,U^nS^jF\rangle_\ell.
\]
The joint spectral resolution of the commuting operators \(U,S\) represents
this as \(\int e^{ink}\lambda^j\,\gamma(dk,d\lambda)\), with
\(0\le\lambda\le1\). Reflection in the row gives the same formula with
\(|j|\) for negative \(j\). Write
\(\lambda=e^{-\sqrt3s}\), including \(s=\infty\) at \(\lambda=0\).
The vertical Fourier coefficients of \(Q_s\) are
\(e^{-\sqrt3s|j|}\). Uniqueness of Fourier coefficients of finite measures
therefore identifies the ordinary increment spectral measure as the mixture
of \(Q_s\)'s with mixing measure \(\gamma\). By
\cref{lem:increment-spectral}, that increment measure is
\(|e^{ik}-1|^2 C'(dk,dl)\). Division on \(k\ne0\) proves
\cref{eq:cauchy-mixture} and the local finiteness assertion.

Since \(Q_s\) has total mass one, the left side of
\cref{eq:weighted-nu} equals \(\int_{0<|k|<\kappa}|k|\,C'(dk,dl)\).
Near the physical origin this is finite by \cref{lem:annulus-mass}, because
\(|k|\le|p|\). The rest consists of finite mass away from the physical
origin, including the alias contribution. This proves the assertion without
any horizontal-band estimate on \(\nu\).
\end{proof}

\subsection{An angular identity with a nonnegative defect}

We seek an angular function whose sixfold rotational average is zero,
but whose unperiodized vertical Poisson average is nonnegative and vanishes precisely
when the normal scale equals the tangential frequency magnitude.  The
fourth harmonic has both properties.  The estimates below control
periodization, aliases, and the passage from horizontal bands to physical
annuli; the dispersion relation is forced only in the scaling limit.

For \(p=(k,l)\ne0\), write
\[
 \chi(k,l)=\cos(4\arg(k+il))
          =1-\frac{8k^2l^2}{(k^2+l^2)^2}.
\]
On \(\mathcal R\), this means the displayed function of the centered
representatives; we only integrate it over small horizontal bands.

\begin{lemma}[Angular Poisson identity]
\label{lem:angular-poisson}
Uniformly for \(s\in[0,\infty]\) and \(0<|k|<\kappa\),
\begin{equation}
 \label{eq:angular-poisson}
 \int\chi(k,l)Q_s(dl)
   =D(k,s)+O(|k|),\qquad
 D(k,s)=\left(\frac{|k|-s}{|k|+s}\right)^2,
\end{equation}
where \(D(k,\infty)=1\).
\end{lemma}

\begin{proof}
Put \(a=|k|\). The bounded analytic function
\[
 z\longmapsto\left(\frac{a+iz}{a-iz}\right)^2
\]
on the upper half-plane has boundary real part \(\chi(a,l)\).
The Poisson formula at \(z=is\) gives, for the unperiodized Cauchy density
\(p_s(l)=s/\{\pi(s^2+l^2)\}\),
\[
 \int_\R\chi(a,l)p_s(l)\,dl
       =\left(\frac{a-s}{a+s}\right)^2.
\]
To bound periodization uniformly, put
\(f_a(l)=1-\chi(a,l)=8a^2l^2/(a^2+l^2)^2\). Then
\(\int_\R f_a(l)\,dl=4\pi a\), and
\(f_a(l)\le 8a^2/l^2\) away from zero. Also
\[
 \sup_{s>0}\sup_{|l|\le P/2}\sum_{n\ne0}p_s(l+nP)<\infty.
\]
For this last bound, split the sum at \(|n|\) comparable to \(s/P\): the
near terms are bounded by \(1/(\pi s)\), and the remaining square tail is
summable. Comparing the integrals of \(f_a\) over the period and over the
line now gives an \(O(a)\) error; the unfolded tail contributes only
\(O(a^2)\). The cases \(s=0\) and \(s=\infty\) follow directly from their
definitions.
\end{proof}

\begin{proposition}[Summable dispersion defect]
\label{prop:dispersion-defect}
For some fixed \(\kappa>0\),
\begin{equation}
 \label{eq:dispersion-defect}
 \int_{0<|k|<\kappa}D(k,s)\,\nu(dk,ds)<\infty.
\end{equation}
There is a finite constant \(K_1\) such that
\begin{equation}
 \label{eq:horizontal-band-mass}
 \nu\{r\le|k|\le2r\}\le K_1
\end{equation}
for all sufficiently small \(r>0\).
\end{proposition}

\begin{proof}
Integrate \cref{eq:angular-poisson} over \(\eps<|k|<\kappa\).
The total error is bounded independently of \(\eps\) by
\cref{eq:weighted-nu}. Replacing the projected measure by the physical
measure in the disk \(|p|<\kappa\) incurs another bounded error: the omitted
physical region and the aliases have finite mass. Thus it suffices to bound
above
\[
 \int_{\eps<|p|<\kappa,\ |k|>\eps}\chi(k,l)\,C(dp).
\]
The integral over the full radial annulus vanishes. Indeed, the sum of the
fourth angular harmonic over rotations by multiples of \(\pi/3\) is zero,
and \(C\) is invariant under those rotations. Write
\(E_\eps=\{|k|\le\eps,\ \eps<|p|<\kappa\}\) for the removed endcap.
Then
\[
 \int_{\eps<|p|<\kappa,\ |k|>\eps}\chi\,dC
 =-\int_{E_\eps}\chi\,dC
 \le\int_{E_\eps\cap\{\chi<0\}}|\chi|\,dC.
\]
Negativity requires \(|l/k|<1+\sqrt2\), so the last integral is confined
to \(\eps<|p|\le K_2\eps\). A fixed number of the annuli in
\cref{eq:annulus-mass} bounds it uniformly. This gives the required upper
bound even if the positive part of the endcap integral is large. Since
\(D\) is nonnegative, monotone convergence gives
\cref{eq:dispersion-defect}.

For \cref{eq:horizontal-band-mass}, on \(r\le|k|\le2r\) the part with
\(s\notin[r/2,4r]\) has \(D\ge1/9\), and hence bounded mass. For
\(s\in[r/2,4r]\), a fixed positive fraction of \(Q_s\) lies in
\(|l|\le2r\). \Cref{eq:cauchy-mixture} bounds that part of \(\nu\)
by a constant times
\(C'\{r\le|k|\le2r,\ |l|\le2r\}\). This is bounded by
\cref{lem:annulus-mass} and the finite alias mass.
\end{proof}

\begin{lemma}[Endcap estimate]
\label{lem:endcap}
For a finite constant \(K_3\),
\begin{equation}
 \label{eq:endcap}
 C\{r<|p|<2r,\ |k|\le\eps\}\le K_3\frac{\eps}{r},
 \qquad 0<\eps<r/10,\quad 0<r<1/4.
\end{equation}
In particular the punctured small disk has no mass on the line \(k=0\).
\end{lemma}

\begin{proof}
Rotate the region by a lattice rotation of angle \(\pi/3\). In the resulting
coordinates \((k',l')\), both \(|k'|\) and \(|l'|\) are comparable to \(r\).
At fixed \(k'\), the original condition \(|k|\le\eps\) restricts \(l'\) to
an interval of length \(O(\eps)\). On this interval the density of every
\(Q_s\), \(0<s<\infty\), is at most \(K/r\): its central Cauchy term is
bounded by \(1/(2\pi|l'|)\), and the periodized tails are uniformly bounded.
The \(s=0\) atom is outside this interval; the uniform law satisfies the same
bound. The relevant horizontal band has bounded \(\nu\) mass by
\cref{eq:horizontal-band-mass}. Since projecting the original local measure
only adds nonnegative alias mass, \cref{eq:cauchy-mixture} proves the estimate
for small \(r\). The remaining compact range of radii is treated identically,
using finiteness of \(\nu\) on horizontal bands away from zero. Finally let
\(\eps\downarrow0\).
\end{proof}

\subsection{Scaling limits of the spectral measure}

For small \(\delta>0\), rescale the physical neighborhood of zero by
\(p\mapsto p/\delta\); denote the resulting measure by \(C_\delta\).
Its restriction to any fixed compact subset of \(\R^2\setminus\{0\}\) is
well defined once \(\delta\) is small. By \cref{lem:annulus-mass}, every
sequence \(\delta\downarrow0\) has a subsequence with a vague limit there.

\begin{lemma}
\label{lem:spectral-subsequences}
Every such vague limit has the form
\begin{equation}
 \label{eq:spectral-subsequence}
 M(dp)=c\,\frac{dp}{|p|^2}
\end{equation}
for a finite \(c\ge0\), initially allowed to depend on the subsequence.
\end{lemma}

\begin{proof}
Fix a compact horizontal band \(a\le|k|/\delta\le b\), where \(0<a<b\).
Rescale the mixing variables to \((k/\delta,s/\delta)\). Their masses are
uniformly bounded by \cref{eq:horizontal-band-mass}. The integral of \(D\)
over the original band tends to zero by \cref{eq:dispersion-defect}.
On the rescaled band the function \(D\) is bounded below away from zero
outside any fixed neighborhood of \(s/\delta=|k|/\delta\), including at
\(s/\delta=\infty\). Therefore every limiting mixing measure is supported
on that graph.

On compact vertical sets, the rescaled periodized Cauchy kernel converges to
the ordinary Cauchy kernel; the folded density error after rescaling is
\(O(\delta)\). Thus, for a nonnegative measure \(\eta\) in the horizontal
coordinate, the limiting physical measure on \(k\ne0\) has the form
\begin{equation}
 \label{eq:limiting-cauchy}
 M(dk,dl)=\frac{|k|}{\pi(k^2+l^2)}\,\eta(dk)\,dl.
\end{equation}
Here aliases introduce no extra term. Their preimages shrink to nonzero
points of the triangular reciprocal torus but, because \(|k|/\delta\ge a\),
exclude those points themselves. Finite measures give vanishing mass to such
shrinking punctured neighborhoods, even when there is an atom at the alias.

\Cref{eq:limiting-cauchy} says that \(|p|^2M(dp)\) is locally
invariant under vertical translations on \(k\ne0\); it does not assume that
\(\eta\) has a density. Apply the same argument after the lattice rotations.
At every nonzero point, at least two of the three tangential coordinates are
nonzero, and the corresponding normal directions are nonparallel. Both
constant directional derivatives of the distribution \(|p|^2M\) vanish in
a neighborhood of that point. It is consequently a constant multiple of
Lebesgue measure there. One can see this by convolving locally with a smooth
approximate identity: the resulting smooth function has both independent
derivatives zero. The constants agree on overlapping neighborhoods, and the
punctured plane is connected. This proves
\cref{eq:spectral-subsequence}, including on the axes. Finiteness and
nonnegativity of \(c\) follow from \cref{lem:annulus-mass} and positivity.
\end{proof}

\subsection{The angular sum fixes the coefficient}

Let \(q\) be any even \(C^\infty\) function with
\begin{equation}
 \label{eq:cutoff-properties}
 0\le q\le1,\qquad q=1\text{ on }[-1/2,1/2],\qquad
 q=0\text{ outside }(-1,1).
\end{equation}
An explicit choice will be given in \cref{sec:normalization}. Define
\[
 w_j(k)=q(2^jk)-q(2^{j+1}k),\qquad
 I_j=\int q(4|p|)\chi(k,l)w_j(k)\,C(dp),\qquad j\ge0.
\]
Each integral is over a compact subset of the punctured torus. No monotonicity
of \(q\) is required.

\begin{proposition}[Angular sum rule]
\label{prop:angular-sum}
The series \(\sum_{j\ge0}|I_j|\) converges. There is a single constant
\(c_*\ge0\) such that every limit in \cref{lem:spectral-subsequences} has
coefficient \(c_*\), and
\begin{equation}
 \label{eq:angular-sum}
 \sum_{j\ge0}I_j=-\frac\pi2 c_*.
\end{equation}
The value is independent of the choice of \(q\) satisfying
\cref{eq:cutoff-properties}.
\end{proposition}

\begin{proof}
The supports of the \(w_j\)'s lie in
\(2^{-j-2}<|k|<2^{-j}\) and have bounded overlap. Replacing the physical
cutoff integral \(I_j\) by
\(\int w_j(k)\chi(k,l)\,C'(dk,dl)\) has errors whose absolute values sum
to a finite number. Indeed, near the physical origin the measures and cutoff
agree, whereas the remaining physical and alias contributions have finite
mass; use the bounded overlap. For all sufficiently large \(j\),
\cref{eq:angular-poisson} then gives
\[
 \int w_j\chi\,dC'=\int w_j(k)D(k,s)\,\nu(dk,ds)+e_j,
 \qquad |e_j|\le K\int |w_j(k)|\,|k|\,\nu(dk,ds).
\]
\Cref{eq:dispersion-defect,eq:weighted-nu} and bounded overlap prove
absolute convergence. The finitely many remaining bands cause no difficulty.

Put \(\eps_N=2^{-N}\). Telescoping, and using \(q(k)=1\) on the support of
\(q(4|p|)\), gives
\[
 \sum_{j=0}^{N-1}I_j
 =\int q(4|p|)\chi(k,l)\bigl[1-q(k/\eps_N)\bigr]C(dp).
\]
Subtract the identically zero radial integral with bracket
\(1-q(|p|/\eps_N)\). The result is
\begin{equation}
 \label{eq:cutoff-difference}
 \sum_{j=0}^{N-1}I_j
 =\int q(4|p|)\chi(k,l)
       \bigl[q(|p|/\eps_N)-q(k/\eps_N)\bigr]C(dp).
\end{equation}
The same right side tends to \(\sum_jI_j\) for arbitrary
\(\eps\downarrow0\). To check this, choose a dyadic \(\eps_N\) comparable
to \(\eps\). The difference of the two horizontal cutoffs is bounded and
supported on a band \(c\eps<|k|<C\eps\). The preceding Poisson comparison
bounds its signed integral by the integrals of \(D\) and \(|k|\) over that
band, plus finite outer and alias mass on a shrinking band excluding
\(k=0\). All three tend to zero. This is also a useful reason for retaining
the integrable defect, rather than only a subsequential limiting formula.

Now take any sequence along which \(C_\eps\) has the limit
\(c\,dp/|p|^2\). After rescaling, the bracket in
\cref{eq:cutoff-difference} vanishes on \(|p|\le1/2\). Its tails are uniform:
on a dyadic annulus \(R<|p|<2R\), \(R>10\), it is supported on \(|k|<1\),
whose rescaled mass is at most \(K/R\) by \cref{lem:endcap}. Summing gives
an \(O(1/R)\) tail bound. Vague convergence may therefore be used in the
whole integral. Polar coordinates yield
\begin{align*}
 \sum_{j\ge0}I_j
 &=c\int_0^{2\pi}\cos(4u)
       \int_0^\infty\bigl[q(r)-q(r|\cos u|)\bigr]\frac{dr}{r}\,du\\
 &=c\int_0^{2\pi}\cos(4u)\log|\cos u|\,du
  =-\frac\pi2c.
\end{align*}
For the middle equality, change variables in the second radial integral,
first with finite cutoffs. The absolute radial integral is bounded by a
constant times \(1+|\log|\cos u||\), which is integrable in \(u\).
For the last equality, use the fourth Fourier coefficient of
\(\log|1+e^{2iu}|\), obtained by expanding from inside the unit circle; the
constant \(-\log2\) in \(\log|\cos u|\) has zero pairing with \(\cos4u\).
The left side is independent of the subsequence, so all coefficients agree.
The final integral also proves independence of \(q\).
\end{proof}

\subsection{Covariance convergence and a reflection null vector}

\begin{theorem}[Universal plane covariance]
\label{thm:plane-covariance}
The rescaled spectral measures converge vaguely on the punctured plane to
\(c_*\,dp/|p|^2\), with \(0<c_*<\infty\). Put \(v=(2\pi)^2c_*\).
For real \(f,g\in C_c^\infty(\R^2)\) of integral zero,
\begin{align}
 \label{eq:plane-covariance}
 \lim_{\delta\downarrow0}\EE[H_\delta(f)H_\delta(g)]
 &=c_*\int_{\R^2}
      \frac{\widehat f(\zeta)\overline{\widehat g(\zeta)}}{|\zeta|^2}\,d\zeta
 \notag\\
 &=\frac{v}{2\pi}\iint_{\R^2\times\R^2}
       f(x)\log\frac1{|x-y|}\,g(y)\,dx\,dy.
\end{align}
The same conclusion holds for tests converging in \(C_c^\infty\) with support
in a common compact set.
\end{theorem}

\begin{proof}
Uniqueness of all vague subsequential limits was proved in
\cref{prop:angular-sum}; relative compactness then gives full vague
convergence. To pass to tests, let
\(A_{\delta,f}(p)=\sum_u a_{\delta,u}(f)e^{ip\cdot u}\) be the transform of
the affine interpolation weights. Their exact total mass is zero. On compact
sets of \(\zeta\),
\(A_{\delta,f}(\delta\zeta)\to\widehat f(-\zeta)\) uniformly. Moreover, for
any fixed integer \(N\ge3\), in a fixed neighborhood of zero,
\begin{equation}
 \label{eq:smearing-transform-bound}
 |A_{\delta,f}(p)|\le K_N
     \min\left\{\frac{|p|}{\delta},
                       (1+|p|/\delta)^{-N}\right\}.
\end{equation}
The first bound follows from zero total mass and
\(\sum_u|a_{\delta,u}|\,|u|=O(\delta^{-1})\). For the second, view
\(a_{\delta,u}\) as the sample at \(u\) of the smooth test convolved with the
compact nodal hat. Each lattice difference gains a factor \(\delta\);
\(N\) summations by parts give the estimate. On the complement of the
neighborhood the same argument gives \(O(\delta^N)\), since at least one
basis difference \(|e^{ip\cdot e_i}-1|\) is bounded below.

Using \cref{lem:annulus-mass}, the squared-transform integral over
\(|p|<\eps\delta\) is \(O(\eps^2)\), and that over
\(R\delta<|p|\) within the neighborhood is \(O(R^{-2N})\). The remaining
part tends to zero by local finiteness and the \(O(\delta^N)\) bound.
On the intermediate annulus, vague convergence and the uniform transform
limit apply. Polarization proves the first equality in
\cref{eq:plane-covariance}. Its second equality is the Fourier representation
of the planar inverse Laplacian on zero-integral tests, whose kernel is
\((2\pi)^{-1}\log(1/|x-y|)\). Height inversion symmetry makes all these
zero-total plane statistics have mean zero. Uniform bounds in
\cref{eq:smearing-transform-bound} prove the assertion for varying tests.

It remains to prove strict positivity. Choose two nonnegative smooth tests
of mass one supported on disjoint disks, and let \(f\) be their difference.
By \cref{lem:buffer-circuits}, with probability bounded below uniformly in
mesh, two nested opposite-sign \(B\) double cuts surround one disk and avoid
the other. The collar contains no pins. The component switch of
\cref{lem:component-switch} then supplies a fair \(W\)-component label which
changes \(H_\delta(f)\) by exactly four when flipped. Conditional on all other
labels its variance is four. Averaging yields a fixed positive lower bound
for \(\EE H_\delta(f)^2\). The covariance limit just proved forces \(c_*>0\).
Its finiteness follows already from the annular mass bound.
\end{proof}

\begin{corollary}[Reflection null vector]
\label{cor:reflection-null}
Let \(\phi\in C_c^\infty(\R^2)\) be supported strictly on one side of a line
parallel to a lattice reflection direction. Choose site rows \(\ell_\delta\)
converging to that line, still separated from the support by a positive
margin. For
\begin{equation}
 \label{eq:laplace-insertion}
 X_\delta=H_\delta(\Delta\phi),
\end{equation}
its reflection norm at \(\ell_\delta\) tends to zero. This holds in every
finite phase augmentation of \cref{prop:reflection-positivity}, and also for
lattice translations and reflected placements converging to the stated
configuration.
\end{corollary}

\begin{proof}
The squared norm is
\(\EE[(\Theta_{\ell_\delta}X_\delta)X_\delta]\). By
\cref{thm:plane-covariance}, it converges to the inverse-Laplacian pairing of
\(\Delta\phi\) and its reflected copy, multiplied by \(v\). That pairing is
\(\int\nabla\phi\cdot\nabla(\phi\circ\theta)\), and is zero because the
supports are disjoint. Reflection positivity makes these squared norms
nonnegative. Auxiliary phases do not change the expectation of a statistic
which does not use them. Converging placements are covered by the varying-test
assertion of \cref{thm:plane-covariance}. The conclusion concerns the reflection
norm; the ordinary variance of \(X_\delta\) need not vanish.
\end{proof}
